% ===== 2.2 Empirical strategy =====

To estimate the effect of antisemitism on views about Israel, the authors start with an OLS regression of the outcome on ``Jews have too much influence'', adding controls step by step (Table~7 of the paper). The coefficient stays stable and even rises as more controls are added. The authors take this as evidence that omitted variables do not push the estimate away from zero, and bounds following \citet{oster2019} suggest that the true effect is even larger. However, three problems remain: further omitted variables, reverse causality, and measurement error, because respondents may hide their antisemitic views \citep[pp.~2352--2355]{binstok2024}.

To address these problems, the authors use an instrumental variable. The idea comes from \citet{voigtlaender2015}: children who grew up under Nazi rule were indoctrinated against Jews in school and in the Hitler Youth, and this shaped their views for the rest of their lives. Figure~\ref{fig:main} confirms that respondents born in the 1930s and 1940s agree more often that Jews have too much influence. Since the State of Israel did not exist at that time, the indoctrination was directed against Jews and not against Israel \citep[pp.~2355--2356]{binstok2024}.

The authors compare respondents born between 1930 and 1945 with those born between 1946 and 1960, while controlling for whether someone was born before 1960. The second and the first stage are:
\begin{equation}\label{eq:main}
  \text{IsraelUnjust}_i = \beta\,\text{JewsInfluence}_i + \delta\,\text{BornBefore1960}_i + \gamma' W_i + \varepsilon_i
\end{equation}
\begin{equation}\label{eq:fs}
  \text{JewsInfluence}_i = \pi\,\text{Born1930--45}_i + \delta'\,\text{BornBefore1960}_i + \gamma'' W_i + u_i
\end{equation}
where $W_i$ contains the controls (in the main specification: gender, education, county share of foreigners and unemployment rate, and federal state), $\beta$ is the effect of interest and $\pi$ measures the strength of the instrument. Standard errors are robust to heteroskedasticity \citep[pp.~2358--2359]{binstok2024}.

The instrument must fulfil two conditions. First, it must be relevant. With a linear control for age, the effect of being born 1930--1945 on antisemitism is too weak; with the control for being born before 1960, it becomes stronger (0.519, standard error 0.118). Second, the instrument must be exogenous. The authors argue that the year of birth cannot be chosen and is therefore exogenously determined. In addition, the exclusion restriction must hold: being born 1930--1945 rather than 1946--1960 may affect views on Israel only through antisemitism. This cannot be tested, and the authors admit that these cohorts also differ in politics, religion and views on foreigners and Islam. They support the assumption in three ways: the direct effect of being born before 1960 is insignificant, adding controls for worldviews increases rather than decreases the estimate, and excluding those born 1940--1945 barely changes the result (0.90 instead of 0.85) \citep[pp.~2356--2360]{binstok2024}.